% =====================================================================
% §7 Discussion — background draft (paragraph-level gesture)
% Target: third open question — external curation of attention
% Length: ~400 words (target band 300–500)
% Generated: 2026-05-29
% Style anchor: paper/sections/06_discussion.tex
%
% Cited from master refs.bib (citeable as-is):
%   - kulveit2025GradualDisempowerment
%   - critchKrueger2020ARCHES
%
% Cited from cluster_C8_algorithmic_curation.bib (AUTO-KEYS — RENAME BEFORE MERGE):
%   - D2025_2          TODO: rename to tkach2025EthicsPersonalization
%   - Lin2026_1        TODO: rename to lin2026AlgorithmicBiasRecommender
%   - Zhao2020_3       TODO: rename to zhao2020InformationCocoons
%
% NB: cluster has 0 cached PDFs; prose is built from bib abstracts/titles
% only. This is deliberate — the section is a forward pointer, not a
% review.
% =====================================================================

\paragraph{External attentional curation.}
The third open question is the most consequential for the model's empirical
domain of validity, and it remains open. Throughout
Sections~\ref{sec:model}--\ref{sec:experiments} an agent's policy over which
neighbours, papers, and observations to attend to is treated as endogenous: it
is selected by expected free energy under the agent's own generative model and
affective utility. In contemporary scientific practice this assumption is
already strained. Citation recommenders, search-engine ranking, social-media
amplification, and journal-level visibility filters increasingly mediate which
findings an individual scientist ever encounters, and the same logic operates a
fortiori in adjacent epistemic communities --- policy analysis, public
discourse around scientific results, the training corpora of large models that
themselves now participate in summarisation and triage. The literature on
algorithmic curation documents that such systems do not merely accelerate
access to a fixed information environment but actively reshape exposure
through popularity bias, preference reinforcement, and the formation of
``information cocoons'' that insulate users from disconfirming evidence
\citep{Lin2026_1, Zhao2020_3}, and that the boundary between helpful
personalisation and manipulation is structurally hard to draw once the
recommender's objective is decoupled from the user's epistemic interest
\citep{D2025_2}.

For the present construction the consequence is precise but unresolved. The
attentional policy $\pi$ no longer factorises through a single agent's
generative model: a curator $\mathcal{C}$ --- algorithmic, institutional, or
both --- intervenes between the latent paradigm field and the per-agent
observation channel, with its own objective that is in general not aligned
with predictive accuracy on the underlying scientific question. The
trust matrix $\trustmat$ then learns precisions over a doubly-mediated
signal, and the resource-flow readouts of Eqs.~\eqref{eq:flow-from-trust}--\eqref{eq:inflow-share}
acquire a third upstream lever beyond $\inflowrate$ and $\trustmat^{(0)}$:
the curator's policy itself. Whether this collapses to a renormalisation of
the existing trust dynamics, or whether it introduces qualitatively new
attractors --- in particular, attractors that survive the corrective force
of cross-surprisal because the disconfirming observations are never sampled ---
is the model's natural next question. It is also the point at which the
construction makes contact with broader concerns about AI systems gradually
displacing human epistemic agency \citep{kulveit2025GradualDisempowerment} and
about the multi-principal coordination problems that arise when curators,
scientists, and institutions optimise different objectives over a shared
information substrate \citep{critchKrueger2020ARCHES}. We flag this as the
open direction the framework most naturally extends into, and the one we
take to be most urgent.
